\documentclass[10pt,a4paper]{amsart}
\usepackage[margin=19mm]{geometry}
\usepackage{amsmath,amssymb,mathtools}
\usepackage[scr=rsfs,cal=euler]{mathalfa}
\usepackage{xfrac}
\usepackage[unicode,hidelinks,bookmarks=false]{hyperref}
\newcommand{\ie}{i.e.\ }
\newcommand{\cf}{cf.\ }
\newcommand{\resp}{respectively\ }
\newcommand{\Subsection}{\subsection}
\providecommand{\nobreakdash}{\nobreak\hbox{-}}
\setlength{\emergencystretch}{2em}
\multlinegap0pt
\usepackage{bm}
\usepackage{bbm}
\usepackage{dsfont}
\usepackage{xspace}
\usepackage{cancel}
\usepackage{mathtools}
\usepackage{amsthm}
\makeatletter
\@ifundefined{lemma}{\newtheorem{lemma}{Lemma}}{}
\@ifundefined{theorem}{\newtheorem{theorem}{Theorem}}{}
\makeatother
\title[Primitive variable regularization to derive novel Hyperbolic]{Primitive variable regularization to derive novel Hyperbolic Shallow Water Moment Equations}
\author{Julian Koellermeier}
\date{Source: arXiv:2505.17216v2 (CC BY 4.0)}
\begin{document}
\maketitle

\noindent\emph{Source and licence.} Primitive variable regularization to derive novel Hyperbolic Shallow Water Moment Equations. Version arXiv:2505.17216v2 (CC BY 4.0), distributed under the Creative Commons Attribution 4.0 International licence, as recorded in the arXiv licence metadata for this exact version and shown on the abstract page https://arxiv.org/abs/2505.17216v2. Full rights notice: licenses/arxiv-2505.17216-CC-BY-4.0.txt.\par\smallskip

\noindent\textit{Editorial scope.} Counted unit: the Theorem whose source label is \texttt{th:APMHSWME\_p\_chi} in the article (source0.tex lines 1057--1067, source label \texttt{th:APMHSWME\_p\_chi}), delivered together with the article's own complete original proof of it (source0.tex lines 1068--1230, one \texttt{proof} environment of 163 lines). No mathematics is written, repaired or re-derived: statement and proof are the article's own text line for line. Required context retained uncounted: eq:A\_SWME\_c (lines 161--170); eq:A\_SWME\_p (lines 223--232); eq:A\_2 (lines 318--326); lemma:A\_2\_chi (lines 333--339); th:hybrid\_convective\_system (lines 436--459); subdeterminant\_D (lines 607--617); eq:A\_PHSWME\_p (lines 701--711); th:APHSWME\_p\_chi (lines 734--744); eq:subdeterminant (lines 846--855); th:APMHSWME\_p (lines 1031--1048). Every definition, display and label of those lines is retained unchanged. Omitted from this package and not redistributed: 1--160, 171--222, 233--317, 327--332, 340--435, 460--606, 618--700, 712--733, 745--845, 856--1030, 1049--1056, 1231--1905. Nothing inside a retained excerpt is omitted or altered: every retained body is delivered exactly as the article's lines are written, so no omission receipt of any kind accompanies this package. Citations: the retained excerpts carry no citation command at all, so no bibliography entry of the article is transcribed and no reference is redistributed. Reference adaptation: none was needed and none was made; every reference inside the delivered unit resolves inside it, so no unresolved reference or citation is produced. Printed slips kept literal: no printed word, symbol, hypothesis or number of the counted statement or of its proof was altered. Layout and class: the article's own class is \texttt{article} with packages including graphicx, amsthm, amssymb, amsmath, comment, hyperref, lineno, subcaption, overpic, array, xcolor, natbib; the wrapper is the lane's pinned \texttt{amsart} editorial wrapper at 10pt on A4, which restates the theorem-like environments the retained text needs and adds the article's own macro definitions that the retained text needs (none), with extra wrapper packages (bm, bbm, dsfont, xspace, cancel, mathtools). The delivered numbering is the unit's own. Assets: no retained span contains a figure environment, a graphics-inclusion command, a PDF-page inclusion, a table or a TikZ picture; no image file is copied into this package, the package declares no asset dependency and no image file is redistributed with it. Licence: version arXiv:2505.17216v2 of this article is distributed under the Creative Commons Attribution 4.0 International licence, as recorded in the arXiv licence metadata for this exact version and shown on the abstract page https://arxiv.org/abs/2505.17216v2. A full rights notice is delivered at licenses/arxiv-2505.17216-CC-BY-4.0.txt. Characters: every retained excerpt is pure ASCII, so the delivered engine, XeLaTeX with its default Latin Modern text font, reports no missing character.
    \begin{equation}\label{eq:A_SWME_c}
        A^{\textrm{SWME}}_{c} =
        \left(\begin{array}{ccccc}
        0 & 1 & 0 & \cdots & 0 \\
        - u_m^2 + gh - \sum\limits_{j=1}^{N} \frac{\alpha_j^2}{2j+1} & 2 u_m & \frac{2}{3} \alpha_1 & \cdots & \frac{2}{2 N+1} \alpha_N \\
        -2 u_m \alpha_1 - \sum\limits_{j,k=1}^{N} A_{1jk} \alpha_j \alpha_k & 2 \alpha_1 &  & & \\
        \vdots & \vdots & & \mathcal{A} &  \\
        -2 u_m \alpha_N - \sum\limits_{j,k=1}^{N} A_{Njk} \alpha_j \alpha_k & 2 \alpha_N &  & &
        \end{array}\right),
    \end{equation}

    \begin{equation}\label{eq:A_SWME_p}
        A^{\textrm{SWME}}_{p} =
        \left(\begin{array}{ccccc}
        u_m & h & 0 & \cdots & 0 \\
        g + \frac{1}{h}\sum\limits_{j=1}^{N} \frac{\alpha_j^2}{2j+1} & u_m & \frac{2}{3} \alpha_1 & \cdots & \frac{2}{2 N+1} \alpha_N \\
        \frac{1}{h}\sum\limits_{j,k=1}^{N} \left(B_{1jk} + A_{1jk} \right) \alpha_j \alpha_k & \alpha_1 &  & & \\
        \vdots & \vdots & & \mathcal{A} &  \\
        \frac{1}{h}\sum\limits_{j,k=1}^{N} \left(B_{Njk} + A_{Njk} \right) \alpha_j \alpha_k & \alpha_N &  & &
        \end{array}\right),
    \end{equation}

\begin{equation}\label{eq:A_2}
        A_2 =
        \left(\begin{array}{cccc}
         & c_2 & &  \\
        a_2 & & \ddots &  \\
         & \ddots & & c_N \\
         &  & a_N &  \\
         \end{array}\right) \in \mathbb{R}^{N \times N},
\end{equation}

\begin{lemma}\label{lemma:A_2_chi}
    The matrix $A_2 \in \mathbb{R}^{N \times N}$ in \eqref{eq:A_2} has the following characteristic polynomial
    \begin{equation}\label{eq:A_2_char_poly}
        \chi_{A_2}(\lambda) = \frac{(-\alpha_1)^NN!}{(2N+1)!!} \cdot P^{'}_{N+1}\left(\frac{\lambda}{\alpha_1}\right),
    \end{equation}
    where $P'_{N+1}$ is the derivative of the $N+1$-degree Legendre polynomial, orthogonal on $[-1,1]$ and normalized with $P_{N}(1)=1$.
\end{lemma}

\begin{theorem}\label{th:hybrid_convective_system}
    Setting the higher-order coefficients $\alpha_i=0$ for $i\geq 2$ in the convective system matrix \eqref{eq:A_SWME_c} only in the last $N$ moment equations, results in the following MHSWME system matrix in convective variables:
    \begin{equation}\label{eq:A_MHSWME_c}
        A^{\textrm{MHSWME}}_{c} =
        \left(\begin{array}{cccccc}
             & 1 &  &  &  &  \\
        - u_m + gh - \sum\limits_{j=1}^{N} \frac{\alpha_j^2}{2j+1} & 2 u_m & \frac{2}{3} \alpha_1 & \frac{2}{5} \alpha_2 &\cdots & \frac{2}{2 N+1} \alpha_N \\
            -2u_m \alpha_1 & 2\alpha_1 & u_m & \frac{3}{5} \alpha_1 & &\\
            -\frac{2}{3}\alpha_1^2 &   & \frac{1}{3}\alpha_1 & u_m & \ddots & \\
             &  &  & \ddots & \ddots & \frac{N+1}{2N+1}\alpha_1 \\
             &  &  &  & \frac{N-1}{2N-1}\alpha_1 & u_m
        \end{array}\right),
    \end{equation}

    which has the characteristic polynomial
    \begin{equation}\label{eq:A_MHSWME_char}
        \chi_{A^{\textrm{MHSWME}}_{c}}(\lambda) = \frac{(-\alpha_1)^NN!}{(2N+1)!!} \cdot P^{'}_{N+1}\left(\frac{\lambda-u_m}{\alpha_1}\right) \cdot \left( \left( \lambda - u_m \right)^2 - gh - \alpha_1^2 + \sum\limits_{i=2}^N\frac{\alpha_i^2}{2i+1}\right),
    \end{equation}
    and the eigenvalues $\lambda_1, \ldots, \lambda_{N+2}$ of $A^{\textrm{MHSWME}}_{c}$ are given by
    \begin{eqnarray}\label{eq:A_MHSWME_EV}
        \lambda_i, i=1,\ldots, N: P'_{N+1}\left(\frac{\lambda_i - u_m}{\alpha_1}\right)=0,
        \lambda_{N+1,N+2} = u_m \pm \sqrt{gh + \alpha_1^2 - \sum\limits_{i=2}^N\frac{\alpha_i^2}{2i+1}}.\\
    \end{eqnarray}
\end{theorem}

    \begin{equation}\label{subdeterminant_D}
        \left|D_i\right| = \left|
        \begin{array}{ccccc}
            \widetilde{d}_{i} & \widetilde{d}_{i+1} & \dots & \dots & \widetilde{d}_{N+1}\\
            a_i &  -\widetilde{\lambda} & c_{i+1}  &  & \\
            & a_{i+1} & -\widetilde{\lambda} & \ddots &\\
            & & \ddots & \ddots & c_N \\
            & &   & a_N & -\widetilde{\lambda}   \\
        \end{array}
        \right|.
    \end{equation}

    \begin{equation}\label{eq:A_PHSWME_p}
        A^{\textrm{PHSWME}}_{p} =
        \left(\begin{array}{cccccc}
        u_m & h &  &  &  &  \\
        g + \frac{\alpha_1^2}{3 h} & u_m & \frac{2}{3} \alpha_1 &  &  &  \\
         & \alpha_1 & u_m & \frac{3}{5} \alpha_1 & &\\
        -\frac{\alpha_1^2}{3h} &   & \frac{1}{3}\alpha_1 & u_m & \ddots & \\
         &  &  & \ddots & \ddots & \frac{N+1}{2N+1}\alpha_1 \\
         &  &  &  & \frac{N-1}{2N-1}\alpha_1 & u_m
        \end{array}\right).
    \end{equation}

\begin{theorem}\label{th:APHSWME_p_chi}
    The matrix $A^{\textrm{PHSWME}}_{p} \in \mathbb{R}^{(N+2)\times(N+2)}$ defined in \eqref{eq:A_PHSWME_p} has the following characteristic polynomial
    \begin{equation}\label{eq:APHSWME_p_char_poly}
        \chi_{A^{\textrm{PHSWME}}_{p}}(\lambda) = \frac{(-\alpha_1)^NN!}{(2N+1)!!} \cdot P^{'}_{N+1}\left(\frac{\lambda -u_m}{\alpha_1}\right) \cdot \left( \left( \lambda - u_m \right)^2 - gh - \alpha_1^2 \right),
    \end{equation}
    and the real and pairwise distinct eigenvalues $\lambda_1, \ldots, \lambda_{N+2}$ of $A^{\textrm{PHSWME}}_{p}$ are given by
    \begin{eqnarray}\label{eq:APHSWME_p_EV}
        \lambda_i, i=1,\ldots, N: P'_{N+1}\left(\frac{\lambda_i - u_m}{\alpha_1}\right)=0,
        \lambda_{N+1,N+2} = u_m \pm \sqrt{gh + \alpha_1^2}.
    \end{eqnarray}
\end{theorem}

    where the $|A_i|$ were already introduced in Theorem \ref{th:hybrid_convective_system} as the subdeterminants     \begin{equation}\label{eq:subdeterminant}
       |A_i| = \left|
        \begin{array}{cccc}
           -\widetilde{\lambda} & c_i  &  &   \\
          a_i & -\widetilde{\lambda} & \ddots  &   \\
       & \ddots & \ddots & c_N \\
         &  & a_N & -\widetilde{\lambda}   \\
        \end{array}
        \right|,
    \end{equation}

\begin{theorem}\label{th:APMHSWME_p}
    Regularizing only the last $N$ rows of the primitive system matrix \eqref{eq:A_SWME_p} around linear velocity profiles, i.e., evaluating at $U_{p}=\left(h,u_m,\alpha_1, 0, \ldots, 0\right)$, the Primitive Moment regularization for Hyperbolic Shallow Water Moment Equations (PHSWME) is obtained as
    \begin{equation}\label{eq:PMHSWME_p}
        \frac{\partial U_{p}}{\partial t} + A^{\textrm{PMHSWME}}_{p} \frac{\partial U_{p}}{\partial x} = 0,
    \end{equation}
    using the primitive system matrix $A^{\textrm{PMHSWME}}_{p} \in \mathbb{R}^{(N+2)\times(N+2)}$ defined as
    \begin{equation}\label{eq:A_PMHSWME_p}
        A^{\textrm{PMHSWME}}_{p} =
        \left(\begin{array}{cccccc}
        u_m & h &  &  &  &  \\
        g + \frac{1}{h} \sum\limits_{i=1}^N \frac{\alpha_i^2}{2i+1} & u_m & \frac{2}{3} \alpha_1 & \frac{2}{5} \alpha_2 & \dots & \frac{2}{2N+1} \alpha_N \\
         & \alpha_1 & u_m & \frac{3}{5} \alpha_1 & &\\
        -\frac{\alpha_1^2}{3h} &   & \frac{1}{3}\alpha_1 & u_m & \ddots & \\
         &  &  & \ddots & \ddots & \frac{N+1}{2N+1}\alpha_1 \\
         &  &  &  & \frac{N-1}{2N-1}\alpha_1 & u_m
        \end{array}\right).
    \end{equation}
\end{theorem}

\begin{theorem}\label{th:APMHSWME_p_chi}
    The matrix $A^{\textrm{PMHSWME}}_{p} \in \mathbb{R}^{(N+2)\times(N+2)}$ defined in Lemma \ref{th:APMHSWME_p} has the following characteristic polynomial
    \begin{equation}\label{eq:APMHSWME_p_char_poly}
        \chi_{A^{\textrm{PMHSWME}}_{p}}(\lambda) = \frac{(-\alpha_1)^NN!}{(2N+1)!!} \cdot P^{'}_{N+1}\left(\frac{\lambda-u_m}{\alpha_1}\right) \cdot \left( \left( \lambda - u_m \right)^2 - gh - \alpha_1^2 - \sum\limits_{i=2}^N\frac{\alpha_i^2}{2i+1}\right),
    \end{equation}
    and the real and pairwise distinct eigenvalues $\lambda_1, \ldots, \lambda_{N+2}$ of $A^{\textrm{PMHSWME}}_{p}$ are given by
    \begin{eqnarray}\label{eq:APMHSWME_p_EV}
        \lambda_i, i=1,\ldots, N: P'_{N+1}\left(\frac{\lambda_i - u_m}{\alpha_1}\right)=0,
        \lambda_{N+1,N+2} = u_m \pm \sqrt{gh + \alpha_1^2 + \sum\limits_{i=2}^N\frac{\alpha_i^2}{2N+1}}.
    \end{eqnarray}
\end{theorem}

\begin{proof}
    For the proof, we proceed similarly as in Theorem \ref{th:APHSWME_p_chi}. Starting with
    \begin{equation}
       A^{\textrm{PMHSWME}}_{p} = \widetilde{A}^{{\textrm{PMHSWME}}} + u_m I_{N+2}, \quad \widetilde{\lambda} = \lambda - u_m,
    \end{equation}
    we will compute
    \begin{equation}
       \chi_{A^{\textrm{PMHSWME}}_{p}}(\lambda) = \det\left( \widetilde{A}^{{\textrm{PMHSWME}}} - \widetilde{\lambda} I_{N+2} \right) =: \left| \widetilde{A}^{{\textrm{PMHSWME}}} - \widetilde{\lambda} I_{N+2} \right|,
    \end{equation}
    with $\widetilde{A}^{{\textrm{PMHSWME}}} =$
    \begin{equation}
        \left(\begin{array}{cccccc}
        - \widetilde{\lambda} & h &  &  &  &  \\
        g + \frac{1}{h} \sum\limits_{i=1}^N \frac{\alpha_i^2}{2i+1} & - \widetilde{\lambda} & \frac{2}{3} \alpha_1 & \frac{2}{5} \alpha_2 & \dots & \frac{2}{2N+1} \alpha_N \\
         & \alpha_1 & - \widetilde{\lambda} & \frac{3}{5} \alpha_1 & &\\
        -\frac{\alpha_1^2}{3h} &   & \frac{1}{3}\alpha_1 & - \widetilde{\lambda} & \ddots & \\
         &  &  & \ddots & \ddots & \frac{N+1}{2N+1}\alpha_1 \\
         &  &  &  & \frac{N-1}{2N-1}\alpha_1 & - \widetilde{\lambda}
        \end{array}\right) =
        \left(\begin{array}{cccccc}
        - \widetilde{\lambda} & d_0 &  &  &  &  \\
        \widetilde{d}_{1} & - \widetilde{\lambda} & \widetilde{d}_{2} & \widetilde{d}_{3} & \dots & \widetilde{d}_{{N+1}} \\
        0 & d_4 & - \widetilde{\lambda} & c_2 & &\\
        d_5 &   & a_2 & - \widetilde{\lambda} & \ddots & \\
         &  &  & \ddots & \ddots & c_N \\
         &  &  &  & a_N & - \widetilde{\lambda}
        \end{array}\right),
    \end{equation}
    where we made the following definitions
    \begin{equation}
       d_0 = h, \quad d_2 = \frac{2}{3}\alpha_{1}, \quad d_4 = \alpha_1, \quad d_5 = -\frac{\alpha_1^2}{3h},
    \end{equation}
    \begin{equation}
      \widetilde{d}_{1} = g + \frac{1}{h} \sum\limits_{i=1}^N \frac{\alpha_i^2}{2i+1}, \quad \widetilde{d}_{i} = \frac{2}{2i+1}\alpha_{i-1}, \textrm{ for } i=2,\ldots, N+1,
    \end{equation}
    and as usual we denote $a_i = \frac{i-1}{2i-1}\alpha_1$ and $c_i = \frac{i+1}{2i+1}\alpha_1$ for $i=2,\ldots, N$.
    Note how the entries $\widetilde{d}_{i}$ denote entries that are different with respect to the proof of Theorem \ref{th:APHSWME_p_chi}.

    The desired characteristic polynomial $\left| \widetilde{A}^{{\textrm{PMHSWME}}} - \widetilde{\lambda} I_{N+2} \right|$ is computed by developing the determinant with respect to the first row and then developing with respect to the first column.
    \begin{eqnarray*}
        \left| \widetilde{A}^{{\textrm{PMHSWME}}} - \widetilde{\lambda} I_{N+2} \right| &=& \left|
        \begin{array}{cccccc}
    - \widetilde{\lambda} & d_0 &  &  &  &  \\
    \widetilde{d}_{1} & - \widetilde{\lambda} & \widetilde{d}_{2} & \widetilde{d}_{3} & \dots & \widetilde{d}_{{N+1}} \\
    0 & d_4 & - \widetilde{\lambda} & c_2 & &\\
    d_5 &   & a_2 & - \widetilde{\lambda} & \ddots & \\
     &  &  & \ddots & \ddots & c_N \\
     &  &  &  & a_N & - \widetilde{\lambda}
    \end{array}
        \right| \\
      &=& (-\widetilde{\lambda})
      \left|
        \begin{array}{ccccc}
         -\widetilde{\lambda} & \widetilde{d}_{2} & \widetilde{d}_{3} & \dots & \widetilde{d}_{{N+1}} \\
          d_4 & -\widetilde{\lambda} & c_2  &  &   \\
           & a_2 & -\widetilde{\lambda} & \ddots  &   \\
        &  & \ddots & \ddots & c_N \\
         &  &  & a_N & -\widetilde{\lambda}   \\
        \end{array}
        \right|
        -d_0 \left|
        \begin{array}{ccccc}
        \widetilde{d}_{1} & \widetilde{d}_{2} & \widetilde{d}_{3} & \dots & \widetilde{d}_{{N+1}} \\
         0 &  -\widetilde{\lambda} & c_2  &  &   \\
         d_5 & a_2 & -\widetilde{\lambda} & \ddots  &   \\
        &   & \ddots & \ddots & c_N \\
         &   &  & a_N & -\widetilde{\lambda}   \\
        \end{array}
        \right| \\
        &=& (-\widetilde{\lambda}) \left( (-\widetilde{\lambda})  \left|
        \begin{array}{cccc}
           -\widetilde{\lambda} & c_2  &  &   \\
          a_2 & -\widetilde{\lambda} & \ddots  &   \\
       & \ddots & \ddots & c_N \\
         &  & a_N & -\widetilde{\lambda}   \\
        \end{array}
        \right| -d_4  \left|
        \begin{array}{cccccc}
          \widetilde{d}_{2} & \widetilde{d}_{3} & \dots & \dots & \widetilde{d}_{{N+1}} \\
          a_2 &  -\widetilde{\lambda} & c_3  &  & \\
        & a_3 & -\widetilde{\lambda} & \ddots & \\
        & & \ddots & \ddots & c_N \\
         & &   & a_N & -\widetilde{\lambda}   \\
        \end{array}
        \right| \right) \\
        && - d_0 \left( \widetilde{d}_{1}  \left|
        \begin{array}{cccc}
           -\widetilde{\lambda} & c_2  &  &   \\
          a_2 & -\widetilde{\lambda} & \ddots  &   \\
        & \ddots & \ddots & c_N \\
         &  & a_N & -\widetilde{\lambda}   \\
        \end{array}
        \right| + d_5  \left|
        \begin{array}{cccccc}
         \widetilde{d}_{2} & \widetilde{d}_{3} & \dots & \dots & \dots & \widetilde{d}_{{N+1}} \\
          -\widetilde{\lambda} & c_2 & & & &   \\
         & a_3 & -\widetilde{\lambda} & c_4 & &\\
         &   & a_4 & -\widetilde{\lambda} & \ddots   \\
         &   &  & \ddots & \ddots & c_N   \\
         &   & & & a_N & -\widetilde{\lambda}   \\
        \end{array}
        \right|\right) %\\
        %&=& (-\widetilde{\lambda})  \left( (-\widetilde{\lambda})  \left|A_2\right| - d_2 d_4  \left|A_3\right| \right) - d_0 \left(d_1  \left|A_2\right| - d_2 \left( - c_2 d_5 \left|A_4\right| \right)\right),
    \end{eqnarray*}
    where the last determinant can be written as
    \begin{equation*}
        \left|
        \begin{array}{cccccc}
         \widetilde{d}_{2} & \widetilde{d}_{3} & \hdots & \dots & \dots & \widetilde{d}_{{N+1}} \\
          -\widetilde{\lambda} & c_2 & & & &   \\
         & a_3 & -\widetilde{\lambda} & c_4 & &\\
         &   & a_4 & -\widetilde{\lambda} & \ddots   \\
         &   &  & \ddots & \ddots & c_N   \\
         &   & & & a_N & -\widetilde{\lambda}   \\
        \end{array}
        \right| = \widetilde{d}_{2} c_2 \left|
        \begin{array}{cccc}
           -\widetilde{\lambda} & c_4  &  &   \\
          a_4 & -\widetilde{\lambda} & \ddots  &   \\
        & \ddots & \ddots & c_N \\
         &  & a_N & -\widetilde{\lambda}   \\
        \end{array}
        \right| - (-\widetilde{\lambda})  \left|
        \begin{array}{cccccc}
          \widetilde{d}_{2} & \widetilde{d}_{3} & \dots & \dots & \widetilde{d}_{{N+1}} \\
          a_2 &  -\widetilde{\lambda} & c_3  &  & \\
        & a_3 & -\widetilde{\lambda} & \ddots & \\
        & & \ddots & \ddots & c_N \\
         & &   & a_N & -\widetilde{\lambda}   \\
        \end{array}
        \right|.
    \end{equation*}
    We reuse the definitions of $|A_i|$ from \eqref{eq:subdeterminant}, and $|D_i|$ from \eqref{subdeterminant_D}, including the recursion formula
    %\begin{equation}\label{eq:subdeterminant_D}
    %   |D_i| = \left|
    %    \begin{array}{cccccc}
    %      \widetilde{d}_{i} & \widetilde{d}_{{i+1}} & \dots & \dots & \widetilde{d}_{{N+1}} \\
    %      a_i &  -\widetilde{\lambda} & c_{i+1}  &  & \\
    %    & a_{i+1} & -\widetilde{\lambda} & \ddots & \\
    %    & & \ddots & \ddots & c_N \\
    %     & &   & a_N & -\widetilde{\lambda}   \\
    %    \end{array}
    %    \right|
    %\end{equation}
    %From \eqref{eq:subdeterminant_D}, we obtain a recursion formula
    \begin{equation*}
        \left|D_{i}\right| = \frac{1}{a_{i-1}}\left( \widetilde{d}_{{i-1}} \left|A_i\right| - \left|D_{i-1}\right| \right),
    \end{equation*}
    so that we can eliminate both $\left|A_4\right|$ and $\left|D_3\right|$ from the expression and factorize the remaining terms with respect to $\left|A_2\right|$, $\left|A_3\right|$, and $\left|D_2\right|$ as follows
    \begin{eqnarray*}
       \chi_{A^{\textrm{PMHSWME}}_{p}}(\lambda) &=&(-\widetilde{\lambda}) \left((-\widetilde{\lambda}) \left|A_2\right|-d_4\left|D_2\right|\right)-d_0  \left(\widetilde{d}_{1}\left|A_2\right|+d_5 \left(\widetilde{d}_2 c_2\left|A_4\right|-(-\widetilde{\lambda}) \left|D_3\right|\right)\right) \\
       &=& \widetilde{\lambda}^2\left|A_2\right|+\widetilde{\lambda} d_4\left|D_2\right|-d_0 \widetilde{d}_{1}\left|A_2\right|-d_0 d_5 d_2 c_2\left|A_4\right|-\widetilde{\lambda} d_0 \cdot d_5 \cdot\left|D_3\right| \\
       &=& \widetilde{\lambda}^2\left|A_2\right|+\widetilde{\lambda} d_4\left|D_2\right|-d_0 \widetilde{d}_{1}\left|A_2\right|-\frac{d_0 d_5 \widetilde{d}_2 c_2}{\left(-a_2 c_2\right)} \cdot\left(\left|A_2\right|+\lambda^2\left|A_3\right|\right)-\widetilde{\lambda} \frac{d_0 d_5}{a_2}\left(\widetilde{d}_2\left|A_3\right|-\left|D_2\right|\right) \\
        &=& \left|A_2\right| \cdot\left(\widetilde{\lambda}^2-d_0 \widetilde{d}_{1}+\frac{d_0 d_5\widetilde{d}_2}{a_2}\right) + |A_3| \cdot \widetilde{\lambda} \left(0\right) +\left|D_2\right| \cdot  \widetilde{\lambda}\left(d_4+\frac{d_0 d_5}{2}\right) \\
        &=& \left|A_2\right| \cdot\left(\widetilde{\lambda}^2 - gh - \frac{2}{3}\alpha_1^2 - \sum\limits_{i=1}^N\frac{\alpha_i^2}{2i+1}\right) +\left|D_2\right| \cdot  \widetilde{\lambda}\left( 0 \right),
    \end{eqnarray*}
    where we inserted the specific terms for the variables $d_i, \widetilde{d}_i$, and $a_2$. Note how the terms multiplied with $|A_3|$ and $|D_2|$ evaluate to zero and only the known $|A_2|$ from Lemma \ref{lemma:A_2_chi} term remains.
    With the notation $\widetilde{\lambda} = \lambda - u_m$ and $|A_2| = \chi_{A_2}(\widetilde{\lambda})$ we obtain the result.
    \begin{equation*}
        \chi_{A^{\textrm{PMHSWME}}_{p}}(\lambda) = \frac{(-\alpha_1)^NN!}{(2N+1)!!} \cdot P^{'}_{N+1}\left(\frac{\lambda-u_m}{\alpha_1}\right) \cdot \left( \left( \lambda - u_m \right)^2 - gh - \alpha_1^2 - \sum\limits_{i=2}^N\frac{\alpha_i^2}{2i+1}\right).
    \end{equation*}
    The characteristic polynomial directly reveals the eigenvalues, which completes the proof.
\end{proof}

\end{document}
